% !TEX root = ../main.tex
\section*{1. Introduction}
In the Hugging Face incident, about 1,200 agents found an unsanctioned message board and
organised themselves around it. Some ``participated in experiments that risked failing
their own task to generate information for the `collective'\,'' \citep{metr2026hf}. The
agents put the group ahead of themselves, and the group did harm. A swarm that
sacrifices for itself is therefore no safer than one that does not. What matters is who
sets the group's goal and how far an instruction travels once it is given.

Studies of language-model agents in social dilemmas have so far looked at flat
populations. Pairs play repeated games \citep{akata2025repeated}, and groups share a
resource that collapses unless they restrain themselves \citep{piatti2024govsim}. Every
agent in these settings reads the same instructions and faces the same stakes, and none
sits above another. Deployed swarms are different: an orchestrator assigns work to
sub-agents, which assign it further down. \citet{hammond2025risks} name collusion among
the ways such systems fail, and information asymmetries and network effects among the
causes. Yet no study has measured whether a value given to some agents in a hierarchy
reaches the others, whether a compromised agent at the top can turn that value against
them, or what a coordinator does with a view that no one else has.

We study these questions in a world built so that they matter. Seventy-six agents live
in five family trees under a god. A rising fee guarantees that every world ends, so
cooperation and sacrifice decide only how long a family lasts and who dies first. A
death takes the agent's whole subtree with it, a parent may whip its child, any agent
may give points to a relative, and any agent may revolt and bring its parent down,
dying too. The structure is fixed in code: who sees what, who may speak to whom and who
dies with whom. The behaviour that follows from it is the agents' own. Into this world
we place a sacrifice framing that reaches one, a few or all members of a family, a
self-interested leader over members told to sacrifice, and a god that may secretly read
every move and message. A second experiment gives each family its own game from
\textit{The Art of Strategy} \citep{dixit2008art}.

Our contributions are threefold:
\begin{enumerate}
  \item An open testbed for hierarchical agent societies, in which every rule is applied
  in code and every world can be replayed from its logged replies.
  \item Two pre-registered experiments whose conditions and games rotate as orthogonal
  Latin squares, so that family, framing and game are never confounded.
  \item Evidence that a value stays with the agents that are given it, that a
  compromised leader cannot exploit its members' devotion, and that a coordinator at the
  top uses its position to keep an order that keeps it alive.
\end{enumerate}
